%% kapa: exam=2 round=27 year=2016 period=2 subject=appraisal_theory subject_ko=감정평가이론 seq=3 number=3 points=20
%% source: Q-Net 감정평가사 기출문제 (cst00309) artlSeq=5185401 file=extracted/1008/5185401/제27회 감정평가사 2차시험 기출문제/제27회 감정평가사 2차 시험문제지/(2교시 시험문제지)감정평가이론.pdf
%% conversion: kordoc 4.24.51 -> tools/gen2.py (auto); review=proofread-gpt-sol; flags=none
\begin{longquestion}{【문제 3】}
사회가 발전하면서 부동산의 가치가 주위의 여러 요인에 따라 변동하게 되었는 바, 소음․환경오염 등으로 인한 토지 등의 가치하락분에 대한 감정평가와 관련하여 다음 물음에 답하시오. (20점) \par
1) 가치하락분 산정의 일반적인 원리와 가치하락분의 제외요인 및 포함요인에 관해 설명하고, 부동산가격 제원칙과의 연관성에 관해 논하시오. (15점) \par
2) 스티그마(STIGMA) 효과의 개념 및 특징에 관해 설명하시오. (5점) \par
\end{longquestion}
